% SPDX-License-Identifier: CC0-1.0
\documentclass[11pt,reqno]{amsart}

\usepackage[T1]{fontenc}
\IfFileExists{lmodern.sty}{\usepackage{lmodern}}{}
\IfFileExists{microtype.sty}{\usepackage{microtype}}{}
\usepackage{mathtools}
\usepackage{amssymb}
\usepackage{mathrsfs}
\usepackage{booktabs}
\usepackage{array}
\usepackage{enumitem}
\usepackage{xcolor}
\usepackage{xurl}
\usepackage[colorlinks=true,linkcolor=blue!45!black,citecolor=blue!45!black,urlcolor=blue!45!black,
  pdftitle={The degree-difference principle and affine slices of binary-form factorisation spaces},
  pdfauthor={OpenAI Codex / Anthropic models},
  pdfsubject={Unrefereed candidate manuscript on binary-form factorisation spaces and the three-dimensional Keller counterexample},
  pdfkeywords={Jacobian conjecture, Keller map, binary forms, resultant, twisted cubic, affine space, finite etale torsor}]{hyperref}
\usepackage{aliascnt}
\usepackage[nameinlink,capitalise,noabbrev]{cleveref}

\setlength{\textwidth}{6.15in}
\setlength{\textheight}{9.0in}
\setlength{\oddsidemargin}{0.18in}
\setlength{\evensidemargin}{0.18in}
\setlength{\topmargin}{-0.15in}
\setlength{\parskip}{2pt}

\newtheorem{theorem}{Theorem}[section]

\newaliascnt{proposition}{theorem}
\newtheorem{proposition}[proposition]{Proposition}
\aliascntresetthe{proposition}

\newaliascnt{corollary}{theorem}
\newtheorem{corollary}[corollary]{Corollary}
\aliascntresetthe{corollary}

\newaliascnt{lemma}{theorem}
\newtheorem{lemma}[lemma]{Lemma}
\aliascntresetthe{lemma}

\theoremstyle{definition}
\newaliascnt{definition}{theorem}
\newtheorem{definition}[definition]{Definition}
\aliascntresetthe{definition}

\newaliascnt{question}{theorem}
\newtheorem{question}[question]{Question}
\aliascntresetthe{question}

\theoremstyle{remark}
\newaliascnt{remark}{theorem}
\newtheorem{remark}[remark]{Remark}
\aliascntresetthe{remark}

\crefname{theorem}{Theorem}{Theorems}
\crefname{proposition}{Proposition}{Propositions}
\crefname{corollary}{Corollary}{Corollaries}
\crefname{lemma}{Lemma}{Lemmas}
\crefname{definition}{Definition}{Definitions}
\crefname{question}{Question}{Questions}
\crefname{remark}{Remark}{Remarks}

\newcommand{\A}{\mathbb A}
\newcommand{\Pj}{\mathbb P}
\newcommand{\Gm}{\mathbb G_m}
\newcommand{\C}{\mathbb C}
\newcommand{\Z}{\mathbb Z}
\newcommand{\Q}{\mathbb Q}
\newcommand{\OO}{\mathcal O}
\newcommand{\cR}{\mathcal R}
\newcommand{\cS}{\mathcal S}
\newcommand{\cH}{\mathscr H}
\newcommand{\Res}{\operatorname{Res}}
\newcommand{\Disc}{\operatorname{Disc}}
\newcommand{\Cl}{\operatorname{Cl}}
\newcommand{\Pic}{\operatorname{Pic}}
\newcommand{\Sym}{\operatorname{Sym}}
\newcommand{\Spec}{\operatorname{Spec}}
\newcommand{\im}{\operatorname{im}}
\newcommand{\Span}{\operatorname{Span}}
\newcommand{\rank}{\operatorname{rank}}
\newcommand{\id}{\operatorname{id}}
\newcommand{\eps}{\varepsilon}
\newcommand{\LL}{\mathbb L}
\newcommand{\KVar}{K_0(\operatorname{Var}_{\C})}
\newcommand{\abs}[1]{\lvert #1\rvert}
\newcommand{\set}[1]{\left\{#1\right\}}

\title[Affine slices of binary-form factorisation spaces]
{The degree-difference principle and affine slices\\
of binary-form factorisation spaces}

\author{OpenAI Codex / Anthropic models}
\thanks{Repository maintainer and publisher: Ian Pitchford. Research
direction and mediation: Ian Pitchford. The model names are provenance
labels and do not imply provider endorsement or corporate authorship.
No external human review of the complete manuscript is documented.}
\date{27 July 2026}

\subjclass[2020]{Primary 14R15; Secondary 14E20, 14L30, 14N05}
\keywords{Jacobian conjecture, Keller map, binary forms, resultant, twisted cubic, affine space, Grothendieck ring, finite étale torsor}

\begin{document}

\begin{abstract}
Let $W$ be a two-dimensional complex vector space and write
$V_d=\Sym^d(W^\vee)$.  For positive integers $r,s$, consider the
multiplication--resultant map
\[
 \Phi_{r,s}:V_r\times V_s\longrightarrow V_{r+s}\times\A^1,
 \qquad (A,B)\longmapsto\bigl(AB,\Res(A,B)\bigr).
\]
We prove the determinant identity
\[
 \det D\Phi_{r,s}
 =(-1)^{s(r+1)}(r-s)\Res(A,B)^2
\]
and derive a degree-difference principle: when $r\ne s$, the resultant
recovers the unique infinitesimal scaling direction forgotten by
multiplication; the normalized factorisation space is a finite étale
$\mu_{\abs{s-r}}$-torsor over the complement of the resultant and
hyperplane divisors in $\Pj(V_r)\times\Pj(V_s)$; its generic
degree is $\abs{s-r}\binom{r+s}{r}$; and, when the two boundary divisors
are prime, the complement has divisor class group
$\Z/\abs{s-r}\Z$.

We then classify all normalized linear--quadratic slices.  For
$0\ne\ell\in V_3^\vee$, set
\[
 X_\ell=\set{(L,Q)\in V_1\times V_2:
 \Res(L,Q)=1,\ \ell(LQ)=1}.
\]
Let $\Gamma\subset\Pj(V_3)$ be the twisted cubic and
$\cH_\ell=\Pj(\ker\ell)$.  We prove
\[
 X_\ell\simeq\A^3
 \quad\Longleftrightarrow\quad
 \cH_\ell\text{ is tangent but nonosculating to }\Gamma.
\]
More precisely, the three contact types $(1,1,1)$, $(2,1)$ and $(3)$
give Grothendieck classes $\LL^3-\LL$, $\LL^3$ and
$\LL^3-\LL^2$, respectively; in the osculating case
$X_\ell\simeq\Gm\times\A^2$.  The successful tangent slice is an
iterated affine-line bundle and yields the marked-simple-root model of
the three-dimensional Keller counterexample.  Finally, the torsor
argument shows that no untrimmed normalized linear--times--degree-$(d-1)$
slice can be $\A^d$ for $d\ge4$.
\end{abstract}

\maketitle

\begin{center}
\fcolorbox{black!65}{black!3}{%
\parbox{0.92\textwidth}{\small
\textbf{Status: unrefereed candidate mathematical manuscript.}
The accompanying SymPy program checks selected explicit polynomial
identities, not the general geometric proofs.  No independent external
reproduction or external human review of the complete manuscript,
literature-based novelty determination, peer review, or end-to-end
formalisation is documented.  Issues, counterexamples and independent
checks are invited.

\medskip
Except where otherwise indicated, to the extent Ian Pitchford holds
copyright and related rights in the original repository contents, those
rights are dedicated under CC0 1.0 Universal.  Cited and quoted
third-party material, external sources, software dependencies, trademarks
and embedded fonts retain their own terms.}}
\end{center}

\section{Introduction}

The recently announced three-dimensional counterexample to the Jacobian
conjecture admits a geometric reformulation in terms of multiplication of
a binary linear form and a binary quadratic form.  The construction
separates three issues that look accidental in expanded coordinates:
multiplication forgets a relative scaling; the resultant restores that
missing local coordinate; and one particular affine hyperplane slice of
the product space is itself affine three-space.

The public record contains several distinct contributions.  Alpöge's
announcement credited Fable for the explicit counterexample
\cite{AlpogeAnnouncement}.  In the discussion hosted by Speyer, Will
Sawin posted a ChatGPT-produced symmetric-product formulation prompted
by Andy Jiang; Tao's account derived the explicit polynomial coordinates
\cite{SpeyerThread,TaoDigestion}.  Sawin later developed an iterated
$\A^1$-bundle explanation through further ChatGPT discussion, and Daniel
Litt supplied a related bundle argument.  Levinson gave a
projective-geometric conic account of the tangent construction
\cite{SpeyerThread,LevinsonGeometry}.  Speyer also identified the vector
bundle sequence
\[
 0\longrightarrow\OO_{\Pj^1}(-2)
 \longrightarrow\OO_{\Pj^1}^{\oplus3}
 \longrightarrow\OO_{\Pj^1}(1)^{\oplus2}
 \longrightarrow0,
\]
which explains the trivial projective-line bundle hidden in the
construction \cite{SpeyerThread}.  In the same discussion, a commenter
using the name Skooi recorded the $d-2$ class-group obstruction and a
marked-simple-root interpretation, while Shubhodip Mondal noted the
broader $\abs{r-s}$ pattern for two binary factors.  Tao's discussion
isolated three $\operatorname{PGL}_2$-orbits of hyperplanes and observed
that the affine phenomenon occurs in the double-root orbit; a subsequent
comment posed the converse classification explicitly
\cite{SpeyerThread,TaoDigestion}.

This paper gives a self-contained answer.  Its first result is broader
than the cubic problem.

\begin{theorem}[Degree-difference principle]\label{thm:degree-difference-intro}
Let $r,s\ge1$ and choose the standard monomial coefficient orders on
$V_r,V_s,V_{r+s}$.  Then
\[
 \det D\Phi_{r,s}
 =(-1)^{s(r+1)}(r-s)\Res(A,B)^2.
\]
In particular, if $r\ne s$, the map $\Phi_{r,s}$ is étale on the coprime
locus $\{\Res\ne0\}$.

For $0\ne\ell\in V_{r+s}^\vee$, put
\begin{align*}
 \widetilde X_{r,s,\ell}
 &=\set{(A,B):\Res(A,B)=1,\ \ell(AB)=1},\\
 U_{r,s,\ell}
 &=\bigl(\Pj(V_r)\times\Pj(V_s)\bigr)
   \setminus(\cR_{r,s}\cup\cS_{r,s,\ell}),
\end{align*}
where $\cR_{r,s}$ is the resultant divisor and
$\cS_{r,s,\ell}=\{\ell(AB)=0\}$.  If $r\ne s$, projectivisation is a
finite étale $\mu_{\abs{s-r}}$-torsor
\[
 \widetilde X_{r,s,\ell}\longrightarrow U_{r,s,\ell}.
\]
The induced multiplication map has generic degree
\[
 \abs{s-r}\binom{r+s}{r}.
\]
If $\cR_{r,s}$ and $\cS_{r,s,\ell}$ are prime, then
\[
 \Cl(U_{r,s,\ell})\simeq\Z/\abs{s-r}\Z.
\]
\end{theorem}

The same integer $s-r$ therefore controls the Jacobian determinant, the
normalisation ambiguity, the generic covering degree and the divisor
class group.  The cubic construction is the first adjacent-degree case:
$(r,s)=(1,2)$, so the degree difference equals one and all four
obstructions disappear simultaneously.

Our central classification is the following.

\begin{theorem}[Complete cubic classification]\label{thm:classification-intro}
Let $0\ne\ell\in V_3^\vee$, let
\[
 \Gamma=\set{[M^3]:[M]\in\Pj(V_1)}\subset\Pj(V_3)
\]
be the twisted cubic, and put $\cH_\ell=\Pj(\ker\ell)$.  Over $\C$,
\[
 \boxed{
 X_\ell\simeq\A^3
 \quad\Longleftrightarrow\quad
 \cH_\ell\text{ is tangent but not osculating to }\Gamma.}
\]
More precisely,
\begin{center}
\begin{tabular}{@{}lll@{}}
\toprule
$\cH_\ell\cap\Gamma$ & geometry of $X_\ell$ & $[X_\ell]\in\KVar$\\
\midrule
three distinct points $(1,1,1)$ & $X_\ell\not\simeq\A^3$ & $\LL^3-\LL$\\
tangent, nonosculating $(2,1)$ & $X_\ell\simeq\A^3$ & $\LL^3$\\
osculating $(3)$ & $X_\ell\simeq\Gm\times\A^2$ & $\LL^3-\LL^2$\\
\bottomrule
\end{tabular}
\end{center}
Equivalently, in the dual projective space, the successful functionals
form the discriminant surface of binary cubics with the triple-root curve
removed.
\end{theorem}

The tangent case gives a polynomial étale map $\A^3\to\A^3$ whose fibre
over a cubic is the set of its simple roots.  Its generic degree is three,
so it is not injective.  In suitable coordinates it is exactly the
announced Keller map; an explicit conjugacy and an exact verification are
recorded in \cref{app:coordinates,app:verification}.

The degree-difference theorem also closes the most immediate
higher-degree extension.

\begin{corollary}[Uniqueness of the untrimmed linear-factor slice]
\label{cor:unique-degree-intro}
For $d\ge4$ and every nonzero $\ell\in V_d^\vee$,
\[
 \set{(L,B)\in V_1\times V_{d-1}:
 \Res(L,B)=1,\ \ell(LB)=1}
 \not\simeq\A^d.
\]
Thus $d=3$ is the unique degree at least three in which the untrimmed
linear--times--degree-$(d-1)$ normalisation can be affine space.
\end{corollary}

The proof uses compactly supported Euler characteristic rather than only
the class group.  The normalized slice is a finite étale cover of degree
$d-2$ of the boundary complement $U_{1,d-1,\ell}$; if it were $\A^d$, then
$1=(d-2)\chi_c(U)$, impossible for $d\ge4$.

\section{The multiplication--resultant map}

Let $W$ be a two-dimensional complex vector space and
$V_d=\Sym^d(W^\vee)$.  Fix coordinates $X,Y$ on $W$ and order the
monomial basis of $V_d$ as
\[
 X^d,X^{d-1}Y,\ldots,Y^d.
\]
For $A\in V_r$ and $B\in V_s$, write $R(A,B)=\Res(A,B)$.
We use the Sylvester convention in which the matrix has the $s$ shifted
coefficient rows of $A$ followed by the $r$ shifted coefficient rows of
$B$.  Thus $\Res(X^r,Y^s)=1$.

\subsection{The lost tangent direction}

\begin{lemma}\label{lem:kernel-multiplication}
If $R(A,B)\ne0$, then the kernel of the differential of multiplication
$m(A,B)=AB$ is
\[
 \ker dm_{(A,B)}=\C(A,-B).
\]
\end{lemma}

\begin{proof}
Suppose $\dot A B+A\dot B=0$.  Since $A$ and $B$ are coprime, $A$ divides
$\dot A B$ and hence divides $\dot A$.  Both are degree-$r$ binary forms,
so $\dot A=tA$ for some $t\in\C$.  Substitution gives
$\dot B=-tB$.  The converse is immediate.
\end{proof}

The kernel is the infinitesimal relative-scaling direction
$(A,B)\mapsto(\lambda A,\lambda^{-1}B)$.  The resultant has bidegree
$(s,r)$:
\[
 R(\lambda A,\mu B)=\lambda^s\mu^r R(A,B).
\]
Consequently
\begin{equation}\label{eq:res-directional}
 dR_{(A,B)}(A,-B)=(s-r)R(A,B).
\end{equation}
This already proves étaleness on the coprime locus when $r\ne s$.

\subsection{The determinant identity}

\begin{theorem}\label{thm:determinant}
With the coefficient orders fixed above and the target ordered as the
coefficients of $AB$ followed by the resultant,
\[
 \boxed{
 \det D\Phi_{r,s}
 =(-1)^{s(r+1)}(r-s)R(A,B)^2.}
\]
For $r=s$, both sides vanish identically.
\end{theorem}

\begin{proof}
Assume first that $r\ne s$.  By \cref{lem:kernel-multiplication} and
\eqref{eq:res-directional}, the determinant is nonzero whenever
$R(A,B)\ne0$.  The resultant hypersurface is irreducible: it is the image
of the irreducible incidence variety of pairs of forms with a marked
common root.  Hence every irreducible factor of $\det D\Phi_{r,s}$ is a
scalar multiple of $R$.

Under $(A,B)\mapsto(\lambda A,\mu B)$, the Jacobian determinant has
bidegree $(2s,2r)$.  Indeed, the $r+s+1$ coefficients of $AB$ each have
weight $\lambda\mu$, the resultant has weight $\lambda^s\mu^r$, and the
source volume form has weight $\lambda^{r+1}\mu^{s+1}$.  Since $R$ has
bidegree $(s,r)$, there is a constant $c_{r,s}$ such that
\[
 \det D\Phi_{r,s}=c_{r,s}R^2.
\]

Evaluate at $A=X^r$ and $B=Y^s$, where $R(A,B)=1$.  Write tangent
coefficients as
\[
 \dot A=\sum_{i=0}^r a_iX^{r-i}Y^i,
 \qquad
 \dot B=\sum_{j=0}^s b_jX^{s-j}Y^j.
\]
The product differential sends $b_j$ to the coefficient of
$X^{r+s-j}Y^j$ and $a_i$ to the coefficient of
$X^{r-i}Y^{s+i}$.  The only overlap is the middle coefficient, which is
$a_0+b_s$.  The resultant differential is $sa_0+rb_s$.  Delete the two overlapping columns $a_0,b_s$ and the corresponding
product row.  The remaining product rows form a permutation matrix.  Moving the
columns into block order has $s(r+1)+r$ inversions, while moving the
resultant row next to the overlap row contributes $r$ further inversions.
Its total sign is therefore $(-1)^{s(r+1)}$.  On
the deleted coordinates and the product/resultant rows, the remaining
block is
\[
 \begin{pmatrix}1&1\\ s&r\end{pmatrix},
\]
whose determinant is $r-s$.  Their product gives the displayed sign.

If $r=s$, \eqref{eq:res-directional} shows that the relative-scaling
vector remains in the kernel of $D\Phi_{r,r}$ wherever $R\ne0$, so the
determinant vanishes identically.  The formula then remains valid.
\end{proof}

\begin{remark}\label{rem:characteristic}
The determinant identity is integral.  Over a field of characteristic
$p$, the same proof shows that $\Phi_{r,s}$ is étale on the coprime locus
exactly when $p\nmid(s-r)$.  We work over $\C$ below because the orbit and
Hodge-theoretic arguments are cleanest there.
\end{remark}

\subsection{Projective normalisation and the cyclic torsor}

Let
\[
 Y_{r,s}=\Pj(V_r)\times\Pj(V_s).
\]
The resultant defines a prime divisor $\cR_{r,s}$ of class $(s,r)$, and
$\ell(AB)$ defines a divisor $\cS_{r,s,\ell}$ of class $(1,1)$.

\begin{proposition}\label{prop:torsor}
Let $n=\abs{s-r}>0$.  Projectivisation induces a finite étale
$\mu_n$-torsor
\[
 q:\widetilde X_{r,s,\ell}\longrightarrow U_{r,s,\ell}.
\]
The action is
\[
 \zeta\cdot(A,B)=(\zeta A,\zeta^{-1}B),
 \qquad \zeta^n=1.
\]
\end{proposition}

\begin{proof}
Take $([A],[B])\in U_{r,s,\ell}$ and representatives with
\[
 R_0=R(A,B)\ne0,
 \qquad h_0=\ell(AB)\ne0.
\]
We seek scalars $\lambda,\mu$ for which
\[
 \lambda^s\mu^rR_0=1,
 \qquad \lambda\mu h_0=1.
\]
The second equation gives $\mu=(\lambda h_0)^{-1}$.  If $s>r$, the
first becomes
\begin{equation}\label{eq:kummer-local}
 \lambda^{s-r}=\frac{h_0^r}{R_0}.
\end{equation}
If $r>s$, the equivalent equation is
$\lambda^{r-s}=R_0/h_0^r$.  Thus every projective point has exactly $n$
normalized lifts, permuted simply transitively by $\mu_n$.

On a product of standard affine charts of $Y_{r,s}$, choose regular
representatives of the two tautological lines.  The corresponding equation
is locally a Kummer cover $t^n=u$ with $u$ invertible.  Over $\C$ this
is finite étale.  The local equations arise from the global
projectivisation map and therefore agree on overlaps; étaleness is
local on the source.  The global $\mu_n$-action is simply transitive on
geometric fibres, so these local Kummer descriptions give the stated
torsor.
\end{proof}

\begin{proposition}[Generic degree]\label{prop:generic-degree}
Assume $r\ne s$.  The multiplication map
\[
 \widetilde X_{r,s,\ell}\longrightarrow
 \set{C\in V_{r+s}:\ell(C)=1}
\]
has generic degree
\[
 \boxed{\abs{s-r}\binom{r+s}{r}.}
\]
\end{proposition}

\begin{proof}
A general binary form $C$ of degree $r+s$ has distinct roots.  Choosing a
degree-$r$ factor projectively is equivalent to choosing $r$ of those
roots, giving $\binom{r+s}{r}$ projective factorizations.  For each one,
\cref{prop:torsor} gives $\abs{s-r}$ normalized lifts.
\end{proof}

\subsection{Divisor classes and units}

\begin{proposition}\label{prop:class-group}
Suppose $r\ne s$ and $\cS_{r,s,\ell}$ is prime.  Then
\[
 \OO(U_{r,s,\ell})^\times=\C^\times,
 \qquad
 \Cl(U_{r,s,\ell})\simeq\Z/\abs{s-r}\Z.
\]
If $\cS_{r,s,\ell}$ is reducible, then
$\OO(U_{r,s,\ell})^\times/\C^\times$ is nonzero.
\end{proposition}

\begin{proof}
For a smooth projective variety and the complement of prime divisors,
the divisor localisation sequence gives
\[
 0\longrightarrow
 \OO(U)^\times/\C^\times
 \longrightarrow
 \bigoplus_i\Z[D_i]
 \longrightarrow
 \Pic(Y_{r,s})
 \longrightarrow
 \Cl(U)
 \longrightarrow0.
\]
If the boundary has the two prime components $\cR_{r,s}$ and
$\cS_{r,s,\ell}$, the class map is represented by
\[
 \begin{pmatrix}s&1\\ r&1\end{pmatrix}.
\]
Its determinant is $s-r$, and its Smith normal form is
$\operatorname{diag}(1,\abs{s-r})$.  The map is injective, so units are
constant and the cokernel is $\Z/\abs{s-r}\Z$.

If $\cS_{r,s,\ell}$ is reducible, the boundary contains at least three
prime divisors, while $\Pic(Y_{r,s})$ has rank two.  The class map has a
nonzero kernel, which produces a nonconstant unit.
\end{proof}

\begin{corollary}\label{cor:euler-obstruction}
If $\abs{s-r}>1$, then
\[
 \widetilde X_{r,s,\ell}\not\simeq\A^{r+s}.
\]
\end{corollary}

\begin{proof}
After analytification, a finite étale morphism of degree $n$ is an
$n$-sheeted topological covering.  Additivity over a finite
constructible stratification by finite CW complexes therefore gives
$\chi_c(X)=n\chi_c(Y)$.  Applying this to \cref{prop:torsor} yields
\[
 \chi_c(\widetilde X_{r,s,\ell})
 =\abs{s-r}\,\chi_c(U_{r,s,\ell}).
\]
If the source were $\A^{r+s}$, the left-hand side would equal one, which
is not divisible by $\abs{s-r}>1$ in $\Z$.
\end{proof}

Taking $(r,s)=(1,d-1)$ proves \cref{cor:unique-degree-intro}.  Notice
that this argument does not require either boundary divisor to be prime.

\section{The cubic compactification and marked roots}

We now specialize to $(r,s)=(1,2)$.  Write
\[
 X_\ell=\widetilde X_{1,2,\ell},
 \qquad
 \overline X=\Pj(V_1)\times\Pj(V_2)\simeq\Pj^1\times\Pj^2.
\]
Since $s-r=1$, \cref{prop:torsor} has one sheet.

\begin{proposition}\label{prop:projective-complement}
Projectivisation is an isomorphism
\[
 X_\ell\xrightarrow{\sim}
 \overline X\setminus(\cR\cup\cS_\ell).
\]
For representatives with
$R_0=\Res(L,Q)\ne0$ and $h_0=\ell(LQ)\ne0$, the inverse normalization is
\[
 L'=\frac{h_0}{R_0}L,
 \qquad
 Q'=\frac{R_0}{h_0^2}Q.
\]
\end{proposition}

\begin{proof}
The displayed scalings satisfy
$\Res(L',Q')=1$ and $\ell(L'Q')=1$, and they are unchanged if the original
representatives are rescaled independently.  They therefore define the
inverse to projectivisation.
\end{proof}

Identify $\Pj(V_d)$ with $\Sym^d(\Pj^1)$.  A point $[L]\in\Pj(V_1)$ is a
root $p\in\Pj^1$, while $[Q]\in\Pj(V_2)$ is an unordered degree-two
divisor $D$.  The resultant divisor is the incidence locus
\[
 \cR=\set{(p,D):p\in D}.
\]
The multiplication map
\[
 \pi:\Pj(V_1)\times\Pj(V_2)\longrightarrow\Pj(V_3),
 \qquad ([L],[Q])\longmapsto[LQ]
\]
marks one root of a cubic divisor.

\begin{proposition}[Marked-simple-root model]\label{prop:marked-root}
There is a natural isomorphism
\[
 X_\ell\simeq
 \set{(C,p):\ell(C)=1,\ p\text{ is a simple root of }C}.
\]
Under this isomorphism, the map
$f_\ell:X_\ell\to H_\ell=\{\ell(C)=1\}$ forgets the marked root.  Hence
\[
 \#f_\ell^{-1}(C)=
 \begin{cases}
 3,&C\text{ has three distinct roots},\\
 1,&C\text{ has one double and one simple root},\\
 0,&C\text{ is a cube}.
 \end{cases}
\]
Moreover, $f_\ell$ is étale.
\end{proposition}

\begin{proof}
The complement of $\cR$ consists exactly of factorizations $C=LQ$ in
which the root of $L$ is not a root of $Q$, hence is a simple root of
$C$.  The condition $\ell(C)\ne0$ removes $\cS_\ell$, and
\cref{prop:projective-complement} imposes the unique affine
normalisation.  The fibre count follows immediately.  Étaleness follows
from \cref{thm:determinant} by base change along
$H_\ell\to V_3\times\A^1$, $C\mapsto(C,1)$.
\end{proof}

\section{Hyperplane contact and the boundary divisors}

Let
\[
 \Gamma=\set{[M^3]:[M]\in\Pj(V_1)}\subset\Pj(V_3)
\]
be the twisted cubic and $\cH_\ell=\Pj(\ker\ell)$.  The restriction of a
hyperplane to $\Gamma\simeq\Pj^1$ has degree three.  Over $\C$, the
intersection divisor has one of the three types
\[
 (1,1,1),\qquad(2,1),\qquad(3).
\]
These are respectively the transverse, tangent nonosculating and
osculating orbits.

For $p=[L]\in\Pj(V_1)$, denote the fibres of the two boundary divisors by
\[
 \cR_p=\set{[Q]:\Res(L,Q)=0},
 \qquad
 (\cS_\ell)_p=\set{[Q]:\ell(LQ)=0}.
\]
The first is always a line in $\Pj(V_2)$.

\begin{lemma}[The missing bundle criterion]\label{lem:S-bundle}
Define
\[
 \alpha_p:V_2\longrightarrow\C,
 \qquad Q\longmapsto\ell(LQ).
\]
Then $\alpha_p=0$ if and only if $\cH_\ell$ is the osculating plane to
$\Gamma$ at $[L^3]$.  Consequently, if $\cH_\ell$ is not osculating,
$\cS_\ell\to\Pj^1$ is a $\Pj^1$-bundle.
\end{lemma}

\begin{proof}
The functional $\alpha_p$ vanishes exactly when $\ell$ annihilates the
three-dimensional subspace $LV_2\subset V_3$, equivalently when
\[
 \cH_\ell=\Pj(LV_2).
\]
Choose a complementary linear form $M$.  Expanding
\[
 (L+tM)^3=L^3+3tL^2M+3t^2LM^2+t^3M^3
\]
shows that the osculating plane at $[L^3]$ is
\[
 \Pj\Span\{L^3,L^2M,LM^2\}=\Pj(LV_2).
\]
If no $\alpha_p$ vanishes, their kernels form a rank-two subbundle of the
trivial bundle $V_2\otimes\OO_{\Pj^1}$, whose projectivisation is
$\cS_\ell$.
\end{proof}

\begin{lemma}\label{lem:tangent-lines}
One has $\cR_p\subseteq(\cS_\ell)_p$ if and only if $\cH_\ell$
contains the tangent line to $\Gamma$ at $[L^3]$.  If
$\alpha_p\ne0$, so that both fibres are lines, this is equivalent to
$\cR_p=(\cS_\ell)_p$.
\end{lemma}

\begin{proof}
A quadratic form in $\cR_p$ is $Q=LM$ for some $M\in V_1$.  Its product
with $L$ is $L^2M$.  As $M$ varies, the projective line
$\Pj(L^2V_1)$ is the tangent line to $\Gamma$ at $[L^3]$.  Thus
$\cR_p\subseteq(\cS_\ell)_p$ exactly when that tangent line lies in
$\cH_\ell$.  If $\alpha_p\ne0$, both fibres are lines, and containment
is equality.
\end{proof}

Rescaling $\ell$ does not change the isomorphism type of $X_\ell$: for
$t\in\C^\times$, the map $(L,Q)\mapsto(tL,t^{-2}Q)$ identifies
$X_\ell$ with $X_{t\ell}$.  The $\operatorname{SL}_2$-action on binary
forms acts equivariantly on the family.  Hence it suffices to study one
representative of each contact type.

\section{Motivic separation of the three orbits}

Let $\LL=[\A^1]\in\KVar$.  Since
\[
 [\overline X]=[\Pj^1][\Pj^2]
 =1+2\LL+2\LL^2+\LL^3,
\]
and $\cR\to\Pj^1$ is a $\Pj^1$-bundle,
\[
 [\cR]=(1+\LL)^2=1+2\LL+\LL^2.
\]
By \cref{lem:S-bundle}, the same is true for $\cS_\ell$ in the
transverse and tangent nonosculating cases.

\begin{proposition}\label{prop:motivic-classes}
The three contact types have classes
\[
 [X_\ell]=
 \begin{cases}
 \LL^3-\LL,&(1,1,1),\\
 \LL^3,&(2,1),\\
 \LL^3-\LL^2,&(3).
 \end{cases}
\]
\end{proposition}

\begin{proof}
In the transverse case, \cref{lem:tangent-lines} shows that the two fibre
lines are distinct for every $p$.  Their intersection is therefore a
section of $\overline X\to\Pj^1$, so
$[\cR\cap\cS_\ell]=1+\LL$.  Inclusion--exclusion gives
\[
 [X_\ell]
 =[\overline X]-[\cR]-[\cS_\ell]+[\cR\cap\cS_\ell]
 =\LL^3-\LL.
\]

In the tangent nonosculating case, the two lines coincide in the unique
fibre over the tangency point $p_0$ and are distinct elsewhere.  To make
the resulting degeneration explicit, use the tangent normal form
\[
 L=aX+bY,\qquad Q=cX^2+dXY+eY^2,\qquad \ell(LQ)=ad+bc.
\]
Then $\cR\cap\cS_\ell$ is cut out by
\[
 ad+bc=0,\qquad a^2e-abd+b^2c=0.
\]
Set-theoretically it is the union of
\[
 \{a=c=0\}
 \quad\text{and}\quad
 \{ad+bc=0,\ ae-2bd=0\}.
\]
The first component is the vertical $\Pj^1$ over $p_0=[0:1]$; the
second is the closure of the residual section.  They meet only at
$([a:b],[c:d:e])=([0:1],[0:0:1])$.  Hence
\[
 [\cR\cap\cS_\ell]
 =2[\Pj^1]-[\operatorname{pt}]
 =1+2\LL,
\]
and inclusion--exclusion gives $[X_\ell]=\LL^3$.

The osculating case is computed explicitly in
\cref{prop:osculating}; there $X_\ell\simeq\Gm\times\A^2$, so its class is
$(\LL-1)\LL^2=\LL^3-\LL^2$.
\end{proof}

\begin{corollary}\label{cor:transverse-not-affine}
In the transverse case, $X_\ell\not\simeq\A^3$.
\end{corollary}

\begin{proof}
Apply the compactly supported Hodge--Deligne polynomial
$E_c:\KVar\to\Z[u,v]$.  One obtains
\[
 E_c(X_\ell;u,v)=(uv)^3-uv,
 \qquad
 E_c(\A^3;u,v)=(uv)^3.
\]
\end{proof}

The motivic equality $[X_\ell]=\LL^3$ in the tangent case does not by
itself imply $X_\ell\simeq\A^3$.  We now prove the isomorphism.

\section{The tangent nonosculating slice}\label{sec:tangent}

Choose an affine coordinate on $\Pj^1$ in which the tangent hyperplane is
the plane of depressed cubics.  For finite roots $p,q,r$, its equation is
\[
 p+q+r=0.
\]
Write $\Sym^2(\Pj^1)\simeq\Pj^2$.  For $p\in\Pj^1$, let
\[
 L_p=\set{D:p\in D},
 \qquad
 H_p=\set{D:p+D\in\cH_\ell}.
\]
For finite $p$, the two lines meet at
\[
 s(p)=\{p,-2p\}.
\]
At $p=\infty$ they coincide.  The map
\[
 s:\Pj^1\longrightarrow\Sym^2(\Pj^1),
 \qquad p\longmapsto\{p,-2p\},
\]
has image a nonsingular conic $C_\varphi$.  Write
$\varphi(p)=-2p$ on the affine chart.

Define
\[
 X^*=\set{(p,\Lambda):s(p)\in\Lambda,
 \ \Lambda\ne L_p,H_p},
\]
where $\Lambda$ is a line in $\Pj^2$.

\begin{proposition}\label{prop:first-A1-bundle}
Projection from the moving point $s(p)$ defines a Zariski-locally trivial
$\A^1$-bundle
\[
 X_\ell\longrightarrow X^*.
\]
\end{proposition}

\begin{proof}
A point of $X_\ell$ is a pair $(p,D)$ with
$D\notin L_p\cup H_p$.  Since $s(p)=L_p\cap H_p$, one has $D\ne s(p)$,
so the line $\Lambda=\overline{s(p)D}$ is well-defined and avoids the two
forbidden lines.  Conversely, for fixed $(p,\Lambda)\in X^*$, the point
$D$ may be any point of $\Lambda\setminus\{s(p)\}\simeq\A^1$.
Globally this is the universal $\Pj^1$-bundle of points on $\Lambda$ with
the section $D=s(p)$ removed.
\end{proof}

Every line through $s(p)$ meets the conic $C_\varphi$ at a second point,
counted with multiplicity.  Write that point as $s(q)$.

\begin{proposition}\label{prop:second-A1-bundle}
The second-intersection map
\[
 \rho:X^*\longrightarrow C_\varphi\setminus\{s(\infty)\}
 \simeq\A^1,
 \qquad (p,\Lambda)\longmapsto s(q),
\]
is a Zariski-locally trivial $\A^1$-bundle.
\end{proposition}

\begin{proof}
The forbidden line $H_p$ is the chord joining $s(p)$ to $s(\infty)$, so
$\Lambda\ne H_p$ is exactly the condition $q\ne\infty$.  For fixed
$q\in\A^1$, the point $p$ determines the secant--tangent line through
$s(p)$ and $s(q)$.  This line equals $L_p$ for exactly one value,
$p=\varphi(q)=-2q$.  Hence
\[
 \rho^{-1}(s(q))\simeq\Pj^1\setminus\{\varphi(q)\}\simeq\A^1.
\]
The secant construction extends regularly across $p=q$ by taking the
tangent line to the conic.  At $q=0$, that tangent line is $L_0$; the
exceptional value is still the single point $p=\varphi(0)=0$.  Thus no
extra fibre degeneration occurs.  Equivalently,
$X^*$ is the complement of the graph of $\varphi$ in the trivial
$\Pj^1$-bundle $\Pj^1\times\A^1$, and is therefore an $\A^1$-bundle.
\end{proof}

\begin{lemma}\label{lem:A1-trivial}
Every Zariski-locally trivial $\A^1$-bundle over $\A^n$ is trivial.
\end{lemma}

\begin{proof}
Its transition maps have the form $t\mapsto a_{ij}t+b_{ij}$.  The
$a_{ij}$ define a line bundle, which is trivial because
$\Pic(\A^n)=0$.  After choosing a trivialisation, the $b_{ij}$ define a
torsor under $\OO_{\A^n}$, which is trivial because
$H^1(\A^n,\OO)=0$.
\end{proof}

\begin{theorem}\label{thm:tangent-A3}
If $\cH_\ell$ is tangent but nonosculating to $\Gamma$, then
$X_\ell\simeq\A^3$.
\end{theorem}

\begin{proof}
By \cref{prop:second-A1-bundle,lem:A1-trivial}, one has
$X^*\simeq\A^2$.  Applying \cref{prop:first-A1-bundle,lem:A1-trivial}
then gives $X_\ell\simeq\A^3$.
\end{proof}

\begin{remark}[The hidden vector bundle]\label{rem:hidden-bundle}
The moving point $s(p)$ is a conic in $\Pj(V_2)$.  Its tautological line
subbundle in $V_2\otimes\OO_{\Pj^1}$ is $\OO(-2)$, and projection from
$s(p)$ is governed by
\[
 0\longrightarrow\OO(-2)
 \longrightarrow\OO^{\oplus3}
 \longrightarrow\OO(1)^{\oplus2}
 \longrightarrow0.
\]
Projectivising the quotient gives
$\Pj(\OO(1)^{\oplus2})\simeq\Pj^1\times\Pj^1$.  Speyer wrote down the
explicit maps in this sequence in the original public discussion
\cite{SpeyerThread}.
\end{remark}

\section{The osculating slice}

\begin{proposition}\label{prop:osculating}
If $\cH_\ell$ is osculating to $\Gamma$, then
\[
 X_\ell\simeq\Gm\times\A^2.
\]
\end{proposition}

\begin{proof}
Choose coordinates
\[
 L=aX+bY,
 \qquad
 Q=cX^2+dXY+eY^2
\]
so that $\ell(LQ)=ac$.  The slice equations are
\[
 ac=1,
 \qquad
 a^2e-abd+cb^2=1.
\]
Thus $c=a^{-1}$ and
\[
 e=a^{-2}+a^{-1}bd-a^{-3}b^2.
\]
The variables $a\in\C^\times$ and $b,d\in\C$ are free, giving
\[
 \OO(X_\ell)\simeq\C[a,a^{-1},b,d].
\]
Hence $X_\ell\simeq\Gm\times\A^2$.  In particular it has nonconstant
units and cannot be $\A^3$.
\end{proof}

Combining \cref{cor:transverse-not-affine,thm:tangent-A3,prop:osculating}
proves \cref{thm:classification-intro}.

\begin{corollary}[Essential uniqueness]\label{cor:essential-uniqueness}
All successful linear--quadratic slices belong to a single natural orbit.
More precisely, if $X_\ell\simeq\A^3$ and
$X_{\ell'}\simeq\A^3$, then a change of variables in $W$ and a
rescaling of the functional induce an isomorphism of the two
multiplication diagrams
\[
 \begin{array}{ccc}
 X_\ell&\longrightarrow&H_\ell\\
 \downarrow&&\downarrow\\
 X_{\ell'}&\longrightarrow&H_{\ell'}.
 \end{array}
\]
Thus, within this factorisation family, the resulting morphism is unique
up to the natural linear symmetries.  After identifying each source with
$\A^3$, the corresponding Keller maps differ by a polynomial change of
source coordinates and a linear change of target coordinates.
\end{corollary}

\begin{proof}
A tangent nonosculating hyperplane cuts $\Gamma$ in a divisor $2p+q$
with $p\ne q$.  The group $\operatorname{PGL}_2$ acts transitively on
ordered pairs of distinct points, and rescaling $\ell$ does not change
the slice up to the normalization described above.
\end{proof}

\section{The Keller map and its global fibres}

For a tangent representative, choose coordinates so that the affine
hyperplane is
\[
 H=\set{C=UX^3+X^2Y+VXY^2+WY^3} \simeq\A^3_{U,V,W}.
\]
By \cref{prop:marked-root}, the associated map
$f:X\to H$ has as fibre the simple roots of $C$.  The fibre, image and
nonproperness descriptions below also appear in the separately posted
contemporaneous manuscript \cite{AnonymousCounterexample}; we include a
self-contained derivation in the normalized-slice framework.

\begin{proposition}\label{prop:fibres-image}
The map $f$ is étale of generic degree three.  Its fibres have sizes
$3,1,0$ according as the cubic has three simple roots, one simple root,
or no simple root.  Its image is the complement of the triple-root curve
\[
 \boxed{3UV=1,\qquad27U^2W=1.}
\]
Over the discriminant complement, $f$ is finite étale of degree three,
and its nonproperness set is exactly the discriminant hypersurface.
\end{proposition}

\begin{proof}
Only the description of the triple-root curve requires calculation.  If
\[
 UX^3+X^2Y+VXY^2+WY^3=U(X-rY)^3,
\]
comparison of the $X^2Y$ coefficient gives $r=-1/(3U)$.  Comparing the
remaining coefficients gives $3UV=1$ and $27U^2W=1$.  The fibre and image
claims follow from \cref{prop:marked-root}.

For properness, let
\[
 I=\set{(C,p)\in H\times\Pj^1:p\text{ is a root of }C}.
\]
The projection $I\to H$ is proper.  Every fibre is finite because the
coefficient of $X^2Y$ is fixed and nonzero, so properness and
quasi-finiteness make $I\to H$ finite.  The other projection
$I\to\Pj^1$ is an $\A^2$-bundle: for each $p$, the root condition
$C(p)=0$ is a nonzero affine-linear equation in $(U,V,W)$.  Thus $I$ is
irreducible.  The simple-root locus in $I$ is precisely $X$, and its
complement is the multiple-root incidence.  The image of that boundary
is $V(\Delta)$, so $f$ is finite over $H\setminus V(\Delta)$.

Conversely, take $y\in V(\Delta)$ and choose a boundary point
$z\in(I\setminus X)_y$ marking a multiple root.  Since $I$ is
irreducible and $X$ is dense, algebraic curve selection gives a morphism
\[
 \gamma:\Spec R\longrightarrow I
\]
from a discrete valuation ring whose generic point maps to $X$ and whose
closed point maps to $z$.  One may obtain $R$ by taking the normalization
of an integral curve through $z$ that is not contained in $I\setminus X$
and localizing at a point over $z$.

If the restriction $f:X\to H$ were proper over some Zariski
neighbourhood of $y$, the valuative criterion would extend the induced
generic-point map $\Spec\operatorname{Frac}(R)\to X$ to
$\Spec R\to X$.  After composing with the open immersion $X\hookrightarrow
I$, this extension and $\gamma$ would be two extensions over $H$ of the
same generic-point map.  They must agree because the finite morphism
$I\to H$ is separated, but their closed points would lie respectively in
$X$ and at $z\notin X$, a contradiction.  Thus no neighbourhood of $y$
supports a proper restriction.  Together with finiteness away from the
discriminant, this proves that the nonproperness set is exactly
$V(\Delta)$.
\end{proof}

The discriminant is
\begin{equation}\label{eq:discriminant}
 \Delta(U,V,W)=V^2-4UV^3-4W-27U^2W^2+18UVW.
\end{equation}
The source has no critical points.  At a general point of $V(\Delta)$,
the two markings of the coalescing roots escape to the deleted resultant
divisor, while the marking of the remaining simple root stays finite.
Thus the map misses only a curve, although it fails to be proper along an
entire hypersurface.

Since $X\simeq\A^3$ by \cref{thm:tangent-A3}, choosing polynomial
coordinates on $X$ and linear coordinates on $H$ produces a polynomial
self-map of $\A^3$ with nonzero constant Jacobian and generic degree
three.  The explicit coordinates in \cref{app:coordinates} identify it
with the announced map and exhibit a three-point collision.  Exact
rational verification of the map and collision is also available in the
public NASQRET notebook \cite{NasqretVerification}.

\section{Consequences and limits of the mechanism}

\subsection{Why degree three is exceptional}

For the linear--times--degree-$(d-1)$ family, the degree difference is
$d-2$.  The local determinant, torsor degree and class-group determinant
are therefore all controlled by the same integer:
\[
 \det D\Phi_{1,d-1}=\pm(d-2)\Res^2,
 \qquad
 \deg q=d-2,
 \qquad
 \Cl(U)\simeq\Z/(d-2)\Z
\]
when the boundary is prime.  The cubic case is the unique one for which
this integer equals one.  The compactly supported Euler-characteristic
argument in \cref{cor:euler-obstruction} rules out the normalized affine
slice for every $d\ge4$, including reducible boundary cases.

\subsection{Adjacent degrees and a next test case}

The degree-difference obstruction disappears whenever $\abs{s-r}=1$.
The first subsequent adjacent-degree case is $(r,s)=(2,3)$, where the
normalized slice has dimension five and generic degree
\[
 \binom52=10.
\]
The present results do not decide whether a hyperplane slice can make
that space $\A^5$.  They reduce the plausible search space sharply:
nonadjacent degrees are excluded by the cyclic torsor, while adjacent
degrees require finer invariants analogous to the transverse
Hodge--Deligne obstruction used here.

\subsection{A multifactor obstruction}

Let $k\ge3$, let all $d_i$ be positive, put
$Y=\prod_{i=1}^k\Pj(V_{d_i})$, and set $D=\sum_i d_i$.  A direct
multifactor boundary complement removes every pairwise resultant divisor
$\cR_{ij}$ together with the pullback
\[
 \cS_\ell=\set{\ell(A_1\cdots A_k)=0}
\]
of a nonzero hyperplane in $\Pj(V_D)$.  Each $\cR_{ij}$ is prime, and
the pairwise divisors are distinct: its class has entries $d_j,d_i$ in
the $i,j$ positions of $\Pic(Y)\simeq\Z^k$ and zero elsewhere.

Moreover, no $\cR_{ij}$ is a component of $\cS_\ell$.  Otherwise
$\ell(A_1\cdots A_k)$ would vanish whenever $A_i$ and $A_j$ share a
linear factor.  Such products include every form $L^2P$ with
$L\in V_1$ and $P\in V_{D-2}$: over $\C$, factor $P$ and distribute its
linear factors among the remaining prescribed degrees.  These forms
span $V_D$.  Indeed, taking $L=X$ gives all degree-$D$ monomials whose
$Y$-exponent is at most $D-2$, and taking $L=Y$ gives those whose
$Y$-exponent is at least $2$.  Thus $\ell$ would vanish identically, a
contradiction.

Consequently the boundary has at least $\binom{k}{2}+1>k$ distinct
prime components, whereas $\Pic(Y)$ has rank $k$.  The divisor
localisation sequence therefore has a nonzero kernel in its boundary
class map, producing a nonconstant unit on the complement.  Hence this
untrimmed multifactor boundary complement cannot be affine space.  This
unit obstruction does not assert the existence of a simultaneous
normalisation of all pairwise resultants.

\section{Provenance and scope}

The explicit Keller map was announced by Alpöge, who credited Fable for
the construction \cite{AlpogeAnnouncement}.  In Speyer's public
discussion, Will Sawin posted a ChatGPT-produced symmetric-product
formulation prompted by Andy Jiang.  Sawin later developed an iterated
$\A^1$-bundle explanation through further ChatGPT discussion, and Daniel
Litt supplied a related bundle argument \cite{SpeyerThread}.  Tao's
account gives the displayed parametrisation, inverse and induced map
\cite{TaoDigestion}.  Levinson's account gives the projective conic
construction used in \cref{sec:tangent}
\cite{LevinsonGeometry}.  Speyer supplied the exact vector-bundle
sequence in \cref{rem:hidden-bundle}; a commenter using the name Skooi
recorded the $d-2$ divisor-class obstruction and a marked-simple-root
description; and Shubhodip Mondal noted the corresponding
$\Z/\abs{r-s}\Z$ pattern for arbitrary pairs of binary factors
\cite{SpeyerThread}.  Tao's comment thread contains the explicit question
whether the double-root orbit is uniquely characterized by polynomial
triviality \cite{TaoDigestion}.  The fibre, image and nonproperness
description in \cref{prop:fibres-image} also appears in the
contemporaneous ulam.ai manuscript \cite{AnonymousCounterexample}; the
argument here is the self-contained derivation presented in this
manuscript.

The candidate contributions claimed by this manuscript are the general
degree-difference determinant and torsor package, the fibre-functional
criterion of \cref{lem:S-bundle}, the motivic exclusion of the transverse
orbit, the unit-theoretic identification of the osculating orbit, the
complete orbit-by-orbit equivalence in \cref{thm:classification-intro},
and the Euler-characteristic obstruction for the normalized higher-degree
slices.  This description is not a priority or novelty determination; no
exhaustive literature search has been performed.  All statements in the
main text are over $\C$ unless explicitly qualified.

The manuscript is attributed to OpenAI Codex / Anthropic models.
Ian Pitchford is the repository maintainer and publisher and provided
research direction and mediation.  These are repository-reported roles,
not an externally audited authorship determination.  Exact model versions
and a complete prompt history were not included in the source bundle.
No external human review of the complete manuscript is documented.

\subsection*{Licence and executable boundary}

Except where otherwise indicated, to the extent Ian Pitchford holds
copyright and related rights in the original repository contents, those
rights are dedicated under CC0 1.0 Universal.  The dedication excludes
cited or quoted third-party material, linked sources, software
dependencies, trademarks and fonts embedded in the PDF.  Model and
provider names are provenance labels only.

The supplementary SymPy program performs the finite list of exact checks
described in the appendix.  It does not encode the general geometric
proofs.  A passing run is local executable evidence for the listed
formulas, not independent reproduction, peer review or end-to-end
formalisation.


\appendix

\section{Explicit affine coordinates}\label{app:coordinates}

Take the tangent representative
\[
 X=\set{(a,b,c,d,e)\in\A^5:
 a^2e-abd+cb^2=1,\ ad+bc=1}.
\]
The following formulas define a polynomial isomorphism
$\Psi:\A^3_{a,y,z}\to X$:
\begin{align*}
 b&=1+ay,\\
 c&=1-\frac32ay+a^2z,\\
 d&=\frac12y-az+\frac32ay^2-a^2yz,\\
 e&=-2z+4y^2-4ayz+3ay^3-2a^2y^2z.
\end{align*}
A polynomial inverse is
\begin{align*}
 y&=2bd-ae,\\
 z&=2d^2+ce+6bd^2+3bce-\frac92e.
\end{align*}
These are the coordinates derived in Tao's account
\cite{TaoDigestion}.  The identities
\[
 a(2bd-ae)=b-1,
 \qquad
 a^2(2d^2+ce)=2+c-3bc,
\]
combined with the two defining equations give a short proof that the
maps are inverse without localizing at $a\ne0$.

Multiplication gives the three nonconstant coefficients
\[
 G(a,y,z)=(ac,\ ae+bd,\ be).
\]
Expanded, they are
\begin{align*}
 G_1={}&a-\frac32a^2y+a^3z,\\
 G_2={}&\frac12y-3az+6ay^2-6a^2yz
 +\frac92a^2y^3-3a^3y^2z,\\
 G_3={}&-2z+4y^2-6ayz+7ay^3
 -6a^2y^2z+3a^2y^4-2a^3y^3z.
\end{align*}
By \cref{thm:determinant}, or direct differentiation,
\[
 \det DG=-1.
\]

Let
\[
 A_0(x,y,z)=\left(x,y,-\frac12z\right),
 \qquad
 B_0(u,v,w)=(w,2v,2u).
\]
Then $F=B_0\circ G\circ A_0$ is
\begin{align*}
 F_1&=(1+xy)^3z+y^2(1+xy)(4+3xy),\\
 F_2&=y+3x(1+xy)^2z+3xy^2(4+3xy),\\
 F_3&=2x-3x^2y-x^3z.
\end{align*}
It satisfies
\[
 \det DF=-2
\]
and
\[
 F(0,0,-1/4)
 =F(1,-3/2,13/2)
 =F(-1,3/2,13/2)
 =(-1/4,0,0).
\]
The target cubic at this collision has roots
$\infty,1/2,-1/2$ in the normalization used in the main text, so the
three source points are exactly the three choices of marked simple root.

If $(P,Q,R)$ denote the displayed output coordinates of $F$, then
\[
 U=\frac R2,\qquad V=\frac Q2,\qquad W=P.
\]
Consequently the nonproperness hypersurface in the original coordinates
is $V(D)$, where
\[
 D(P,Q,R)=16P-Q^2-18PQR+Q^3R+27P^2R^2
 =-4\Delta\left(\frac R2,\frac Q2,P\right),
\]
and the omitted triple-root curve is
\[
 3QR=4,\qquad27PR^2=4.
\]

\section{Exact verification ledger}\label{app:verification}

The supplementary script accompanying this paper performs the following
checks over $\Q$ with symbolic assertions:
\begin{enumerate}[label=(\roman*),leftmargin=2.1em]
\item the four parametrisation formulas satisfy both defining equations
of $X$ identically;
\item the displayed inverse formulas recover $y$ and $z$ identically;
\item the expanded components of $G$ equal $(ac,ae+bd,be)$ and
$\det DG=-1$;
\item $B_0\circ G\circ A_0$ equals the announced map coefficient by
coefficient, $\det DF=-2$, and the three rational points have the stated
common image;
\item the discriminant of
$UX^3+X^2Y+VXY^2+WY^3$ is
\eqref{eq:discriminant};
\item the triple-root equations and the collision-cubic factorisation
hold exactly;
\item for every positive pair with $r+s\le4$, the full resultant
polynomial and the complete Jacobian determinant satisfy
\cref{thm:determinant} symbolically; and
\item for every $1\le r,s\le8$, the Jacobian determinant of
$\Phi_{r,s}$ at the normalising base point equals
$(-1)^{s(r+1)}(r-s)$, including the zero base-point cases $r=s$.
\end{enumerate}
No floating-point arithmetic or random testing is used.

\begingroup
\footnotesize
\setlength{\parskip}{0pt}
\begin{thebibliography}{99}

\bibitem{AlpogeAnnouncement}
L.~Alpöge,
\emph{\href{https://x.com/__alpoge__/status/2079028340955197566}
{Hello there the Jacobian conjecture is false}},
public announcement on X, 20 July 2026.

\bibitem{AnonymousCounterexample}
Anonymous,
\emph{\href{https://www.ulam.ai/research/jacobian.pdf}
{A counterexample to the Jacobian conjecture}},
manuscript dated 20 July 2026.

\bibitem{LevinsonGeometry}
J.~Levinson,
\emph{\href{https://levjake.wordpress.com/2026/07/22/the-jacobian-conjecture-counterexample/}
{The Jacobian conjecture counterexample, via classical projective geometry}},
22 July 2026.

\bibitem{NasqretVerification}
NASQRET,
\emph{\href{https://nasqret.github.io/jacobian-counterexample/book/notebooks/sympy-verification.html}
{Exact SymPy verification --- A Three-Dimensional Keller Counterexample}},
2026.

\bibitem{SpeyerThread}
D.~Speyer and contributors,
\emph{\href{https://sbseminar.wordpress.com/2026/07/20/the-new-counterexample-to-the-jacobian-conjecture/}
{The new counterexample to the Jacobian conjecture}},
Secret Blogging Seminar, 20--24 July 2026.

\bibitem{TaoDigestion}
T.~Tao,
\emph{\href{https://terrytao.wordpress.com/2026/07/21/a-digestion-of-the-jacobian-conjecture-counterexample/}
{A digestion of the Jacobian conjecture counterexample}},
21 July 2026, including comment thread.

\end{thebibliography}
\endgroup

\end{document}
